\begin{proof}
    For each $n \in \mathbb{N}$, let $D_n: \Core[n] \to \Core[n]$ denote the unique prop functor such that $D_n \circ i = \ddagger \circ i$.
    The proof follows by induction on $n$.
    \begin{itemize}
    \item \textbf{Base Case}.
          Let $g \in \Sigma_{\Core}^0$.
          Then by definition, $D_0( g ) = g^\dagger$.
          Since $g$ was arbitrary, then $D_0 \circ i = \dagger \circ i$.
          Then by uniqueness, $D_0( U ) = U^\dagger$ for all $U \in \Mor( \Core[0] )$.
          Then by reflexivity, $D_0( U ) \approx_E U^\dagger$ for all $U \in \Mor( \Core[n] )$.
    \item \textbf{Inductive Hypothesis}.
          For some $n \in \mathbb{N}$, if $U \in \Mor( \Core[n] ) $, then $D_0( U ) \approx_E U^\dagger$.
    \item \textbf{Inductive Step}.
          Assume that for an $n \in \mathbb{N}$, if $U \in \Mor( \Core[n] )$, then $D_n( U ) \approx_E U^\dagger$.
          Let $g \in \Sigma_{\Core}^{n+1}$.
          If $g \in \Sigma_{\Core}^n$ as well, then $D_{n+1}( U ) = D_n( U ) \approx_E g$ by the inductive hypothesis.
          Assume instead that $g \not \in \Sigma_{\Core}^n$.
          Then there exists some $U \in \Mor( \Core[n] )$ such that $g = \scalebox{0.65}{\circimg{tikz-cache/dagger__term1.pdf}} \odot D_n( U )$.
          Then by the inductive hypothesis, $D_n( U ) \approx_E U^\dagger$.
          Then the following derivation holds.
          \begin{align*}
              \scalebox{0.65}{\circimg{tikz-cache/completeness__ddag_1.pdf}}
              \;\;\;\overunderset{\HQCCtrlCong}{I.H}{\approx}\;\;\;&
              \scalebox{0.65}{\circimg{tikz-cache/completeness__ddag_2.pdf}}
              \overset{\text{Thm(\ref{Thm:HQCInv})}}{\approx}
              \scalebox{0.65}{\circimg{tikz-cache/completeness__ddag_3.pdf}}
              \overunderset{\HqcCtrlComp}{I.H}{\approx}\;\;\;
              \scalebox{0.65}{\circimg{tikz-cache/completeness__ddag_4.pdf}}
              \\
              \overunderset{\HQCCtrlCong}{\text{Thm(\ref{Thm:HQCInv})}}{\approx}&
              \scalebox{0.65}{\circimg{tikz-cache/completeness__ddag_5.pdf}}
              \overset{\HqcCtrlId}{\approx}
              \scalebox{0.65}{\circimg{tikz-cache/completeness__ddag_6.pdf}}
          \end{align*}
          In either case, $D_{n+1}( g ) \approx_E g^\dagger$.
          Since $g$ was arbitrary, then $D_{n+1} \circ i = \dagger \circ i$.
          Then by uniqueness, $D_{n+1}( U ) \approx_E U^\dagger$ for all $U \in \Mor( \Core[n] )$.
          Then the inductive step holds.
    \end{itemize}
    Then by the principle of induction, for each $n \in \mathbb{N}$, if $U \in \Mor( \Core[n] )$, then $D_n( U ) \approx_E U^\dagger$.
    Then let $U \in \Mor( \Core )$.
    There there exists some $n \in \mathbb{N}$ such that $U \in \Mor( \Core[n] )$.
    Then $U^\ddagger = D_n( U ) \approx_E U^\dagger$.
    Since $U$ was arbitrary, then $U^\ddagger \approx_E U^\dagger$ for all $U \in \Mor( \Core )$.
\end{proof}
